%!TEX Root = ../main.tex
\chapter{Modules}
\label{chap:modules}
There are two key reasons why we want to study modules
in commutative algebra.
\begin{itemize}
    \item Modules is an analogue of vector spaces for rings.
    Thus, studying modules will allow us to do linear algebras
    over a ring.
    \item Modules unify various structures in commutative algebra,
    including ideals $I$, quotient rings $R/I$,
    and their direct sum / direct product,
    allowing us to treat those with an equal footing.
\end{itemize}
This chapter introduces basic notions and properties of modules.
\section{Modules}
\subsection{Definition and Basic Properties}
Given a ring $R$, a \vocab{module} $M$ is an abelian group $(M,+)$
with mutliplication by elements in $R$:
a map $R\times M\to M$ satisfying
such that multiplication is compatible with ring operations:
\begin{itemize}
    \item $a(m_1+m_2) = am_1+am_2$ for all $a\in R$ and 
    $m_1,m_2\in M$.
    \item $(a+b)\cdot m = am + bm$ for all $a,b\in R$ and $m\in M$.
    \item $a\cdot (bm) = (ab)m$ for all $a,b\in R$
    and $m\in M$.
\end{itemize}
\begin{example}
\begin{itemize}
    \item When $R=k$, a field, $R$-module is a $k$-vector space.
    \item When $R=\mathbb Z$, a $\mathbb Z$-module is just
	an abelian group.
    \item When $R=\mathbb C[x]$, an $\mathbb C[x]$-module
	is a $\mathbb C$-vector space associated with
	a linear map representing a multiplication by $x$.
\end{itemize}
\end{example}
If there is a subset $N\subset M$ preserved by
addition and scalar multiplication, then
$N$ is an $R$-\vocab{submodule} of $M$.

We can also define an $R$-\vocab{module homomorphism}
to be a map between two $R$-modules $\phi : M\to N$
that is a group homomorphism and
$\phi(am) = a\phi(m)$.

The category of $R$-module is denoted $\Mod R$.
\notation{$\Mod R$}{the category of $R$-modules}
%
%
%
%
\subsection{Constructions of Modules}
We have some constructions of modules.

\paragraph{Quotient}
For any submodule $N\subset M$,
one can define the \vocab{quotient module}
$M/N$ \notation{$M/N$}{quotient module}
by defining operation on cosets $\{N+a : a\in M\}$
as usual.
We also have the following universal property:

For any $R$-module $P$ and a map $\phi : M\to P$,
such that $\phi(N)=0$, then there exists unique
ring homomorphism $\overline\phi : M/N\to P$
such that $\phi=\overline\phi\circ\pi$.

\paragraph{Kernel, Image, Cokernel}
Given $\phi : M\to N$, then 
\notation{$\Ker\phi$}{kernel of module morphism}
\notation{$\Im\phi$}{image}
\notation{$\Coker\phi$}{cokernel}
$$\begin{array}{rrcll}
    \text{\vocab{kernel}}\quad &
    \Ker\phi &:=& \{m\in M : \phi(m)=0\}
    &\ \subseteq\ M\\[4pt]
    \text{\vocab{image}}\quad &
    \Im\phi &:=& \{\phi(m) : m\in M\} 
    &\ \subseteq\ N\\[4pt]
    \text{\vocab{cokernel}}\quad &
    \Coker\phi &:=& N/\Im\phi
\end{array}$$
are all modules.
The following isomorphism theorems follow similarly
from ones in group theory:
\begin{itemize}
    \item $M/\Ker\phi\simeq\Im\phi$
    \item Given a chain of modules $L\subseteq M\subseteq N$,
    we have $(N/L)/(M/L) \simeq N/M$.
\end{itemize}
\paragraph{Hom}
\notation{$\Hom_R(M,N)$}{$\Hom$ module}
Given two modules $M,N$, $\Hom_R(M,N)$ is also a module.
The operation is defined by
\begin{itemize}
    \item $(\phi_1 + \phi_2)(m) = \phi_1(m) + \phi_2(m)$ 
    \item $(a\phi)(m) =a\phi(m)$.
\end{itemize}
\notation{$\End_R(M)$}{endomorphism module}
When, $M=N$, there is another notation $\End_R(M)
:= \Hom_R(M,M)$, which is a noncommutative ring
whose multiplication is given by composition.
\paragraph{Submodules Generated by Ideal}
Given an $R$-module $M$ and $I\subset R$,
we can define a module
\notation{$IM$}{submodule generated by ideal}
$$IM := \left\{\sum_{i=1}^k a_im_i : a_i\in I, m_i\in M\right\}.$$
\paragraph{Annihilator}
We have the \vocab{annihilator} of an element:
\notation{$\Ann(m)$}{annihilator of an element}
$$\Ann(m) := \{a\in R : am=0\}$$
and of the entire module
\notation{$\Ann(M)$}{annihilator of the module}
$$\Ann(M) := \{a\in R : am=0\text{ for all }m\in M\}.$$

\paragraph{Restriction of Scalar}
Let $\kappa : R\to S$ be a ring
homomorphism. Let $N$ be a $S$-module,
then we can define the \vocab{restriction of scalar},
which is an $R$-module $_R N$
\notation{$_R N$}{restriction of scalar}
(or $\Res_R(N)$) by using the same abelian group $N$
and define $a\cdot n = \kappa(a)\cdot n$.
This is a functor from $\Mod S$ to $\Mod R$.
\paragraph{Direct sum and direct product}
Given a collection of modules $(M_\alpha)_{\alpha\in T}$, 
the \vocab{direct product} (also known as \vocab[product of modules]{product})
is defined by
\notation{$\prod_{\alpha\in T}M_\alpha$}{product of modules}
$$\prod_{\alpha\in T} M_\alpha
:= \{(m_\alpha)_{\alpha\in T}\},$$
and all operations are done independently on each coordinate.
(If it is a finite collection, we write
$M_1\times M_2 \times \dots\times M_n$.)
\notation{$M_1\times\dots\times M_n$}
{product of modules}
The universal property is similar to universal property
of products in \Cref{sec:universal}:
if $\tau_\alpha : P \to M_\alpha$, then
there exists unique $\tau : P\to \prod M_\alpha$.


The \vocab{direct sum} (also known as \vocab{coproduct}) 
is defined as
\notation{$\bigoplus_{\alpha\in T}M_\alpha$}{direct sum of modules}
\begin{align*}
\bigoplus_{\alpha\in T} M_\alpha
= \Big\{(m_\alpha)_{\alpha\in T} : \text{
only finitely many } m_\alpha \text{ are nonzero.}\Big\}
\end{align*}
(If it is a finite collection, we write
$M_1\oplus M_2 \oplus \dots\oplus M_n$.)
\notation{$M_1\oplus\dots\oplus M_n$}
{direct sum of modules}
Let $i_\alpha : M_\alpha\to\bigoplus M_\alpha$
be the inclusion map. 
Then, we have the universal property:
for any $\phi_\alpha : M_\alpha\to P$,
then there exists unique $\phi : \bigoplus M_\alpha\to P$
such that $\phi_\alpha = \phi\circ i_\alpha$.
This is exactly the same as the universal property of product,
but all arrows are reversed.

Notice that finite direct sum is the same as finite product,
but infinite direct sums and products are not the same.


A \vocab[free module]{free} $R$-module is the direct sum of
copies of $R$. We write it as $R^{\oplus n}$
for integer $n\geq 0$ 
or $R^{\oplus T}$ for some set $T$.
\notation{$R^{\oplus n}$}{free module}

\paragraph{Generators}
An $R$-module $M$ is a
\begin{itemize}
    \item \vocab[finitely-generated module]{finitely-generated}
    (or just \vocab[finite module]{finite}) if 
    there exists a surjective map $\phi : R^{\oplus n}\to M$
    for some integer $n\geq 0$.
    \item \vocab[finitely-presented module]{finitely-presented} if 
    the $\phi$ above exists and $\Ker\phi$ is finitely-generated.
    \item \vocab[cyclic module]{cyclic} if the $\phi$ above exists with $n=1$
    (i.e., $M=R/I$ for some ideal $I$).
\end{itemize}
%
%
%
%
%
\subsection{Modules Over PID}
We recall the following theorem from undergraduate algebra.
\begin{theorem}[PID Structure Theorem]
    \index{PID Structure Theorem}
    \label{thm:pid_structure}
    Let $R$ be a PID, then every \textbf{finitely-generated} $R$-module $M$ 
    is a direct sum of cyclic modules. 
    In particular, there exists $r\geq 0$
    and $a_1,\dots,a_n\in R$ such that
    $$M\simeq R^{\oplus r} \oplus \bigoplus_{i=1}^n 
    R/(a_i).$$
\end{theorem}
\begin{proof}
    At least in the case when $R$ is a Euclidean domain,
    this was done in undergraduate algebra
    \cite[Theorem 14.7.3]{artin}.
    The process of obtaining Smith normal form
    can be modified slightly to work on any PID.
    In particular, one must replace the use of Euclidean division
    to using the fact that for all $a,b\in R$,
    there exists $x,y\in R$ such that $ax+by=\gcd(a,b)$.
\end{proof}
\begin{exercise}
Prove that $\QQ$ is not a finitely-generated $\ZZ$-module.
\end{exercise}
%
%
%
%
%
\section{Cayley-Hamilton Theorem and Consequences}
\label{sec:cayley}
\subsection{Cayley-Hamilton Theorem}
\begin{theorem}[Cayley-Hamilton Theorem]
    \index{Cayley-Hamilton Theorem}
    \label{thm:cayley_hamilton}
    Let $R$ be a ring. Let $A=\Mat_{n\times n}(R) = \big(a_{ij}\big)_{1\leq i,j\leq n}$.
    We define the \vocab{characteristic polynomial}
    by 
    $$p_A(t) = \det(tI_n-A)=\sum_{k=0}^n c_k t^k.$$ 
    Then, $p_A(A) = 0$.
\end{theorem}
\begin{proof}
    The proof is in three steps.

    \textbf{Step 1. $R=\CC$ and $A$ has distinct eigenvalues.}
    Let the eigenvalues be $\lambda_1,\dots,\lambda_n$
    and the corresponding eigenvectors be $v_1,\dots,v_n$.
    Then, since
    $$p_A(t) = (t-\lambda_1)\dots (t-\lambda_n),$$
    we get that $p_A(A)v_j = 0$ for all $j$ because the factor
    $A-\lambda_j$ alone is sufficient to annihilate $v_j$.
    Since $v_1,\dots,v_n$ spans $\CC^n$,
    it follows that $p_A(A)v=0$ for all $v\in\CC^n$,
    so $p_A(A)=0$.
    
    \textbf{Step 2. $R=\mathbb Z[x_{11},x_{12},
    \dots x_{nn}]$ and $A$ is of the following form}
    $$A = \begin{pmatrix}
	x_{11} & \dots & x_{1n} \\
	\vdots & \ddots & \vdots \\
	x_{n1} & \dots & x_{nn}
    \end{pmatrix}.$$
    To prove this, simply replace $x_{ij}$
    with a complex number so that $x_{ij}$ are
    transcendentally independent.

    \textbf{Step 3. General $R$.}
    From the above case, substitute 
    $x_{ij}$ with $a_{ij}\in R$, and since 
    $p_A(A)$ vanish as polynomial ring $\ZZ[x_{11},
    \dots,x_{nn}]$, it must still vanish when plugging in
    $x_{ij}$ with $a_{ij}$.
\end{proof}
The above proof shows a more general principle of 
\vocab{permanence of identities}:
if one has identity that is true over $\CC$,
one can carry through the above steps to show
that such identity hold for any ring.

We now explain some consequence of this theorem.
\begin{corollary}[Cayley-Hamilton; module version]
    \label{cor:cayley_hamilton}
    \index{Cayley-Hamilton Theorem}
    Let $M$ be a finitely-generated $R$-module.
    Let $\phi : M\to M$ be an endormorphism
    such that $\phi(M) \subset IM$ for some ideal $I\subset R$.
    Then,
    $$\phi^n + c_{n-1}\phi^{n-1} + \dots + c_1\phi + c_0 
    = 0,$$
    for some $c_0,\dots,c_{n-1}\in I$.
\end{corollary}
\begin{proof}
    View $M$ as a quotient of $R^{\oplus n}$ for some $n$.
    Let $e_1,\dots,e_n$ be a basis of $R^{\oplus n}$,
    so their reductions $\overline e_1, \dots, \overline e_n$
    generates $M$.

    We claim that there exists $a_{ij}\in I$ such that
    $\phi(\overline e_i) = \sum_{j=1}^n a_{ij} \overline e_{ij}$.
    To see this, we note that since $\phi(\overline m_i)
    \in IM$, it is a sum of elements in form $im$
    for $i\in I$ and $m\in M$.
    However, $m$ is an $R$-linear combination of $\overline e_1,\dots,\overline e_n$,
    so $im$ (and hence $\phi(\overline m_i)$)
    is an $I$-linear combination of $\overline e_1,\dots,\overline e_n$.

    The elements $(a_{ij})_{1\leq i,j\leq n}$ form a matrix 
    $A\in \Mat_{n\times n} (R)$.
    Thus, by Cayley-Hamilton theorem (\Cref{thm:cayley_hamilton}), we get that
    there exists $c_0,\dots,c_{n-1}$ such that
    $$A^n + c_{n-1}A^{n-1} + \dots + c_1A + c_0 = 0.$$
    Reducing this equation into $M$ gives
    $$\phi^n + c_{n-1}\phi^{n-1} + \dots + c_1\phi + c_0=0.$$
    Note that $c_0,c_1,\dots,c_{n-1}\in I$
    because $c_0,\dots,c_{n-1}$ come from characteristic polynomial of $A$, and each entry of $A$ is in $I$. 
    Hence, we are done.
\end{proof}
\subsection{Nakayama's Lemma}
We will now use \Cref{cor:cayley_hamilton}
to deduce an important result called \vocab{Nakayama's lemma}
that will appear sporadically throughout the text.
(If you don't believe me, Ctrl+F the word ``Nakayama''
in this file.)
\begin{lemma}[Nakayama's Lemma]
    \index{Nakayama's lemma}
    Suppose that $M$ is a finitely-generated $R$-module
    such that $IM=M$, then
    there exists $i\in I$ such that $im=m$
    for all $m\in M$.

    (Mnemonic: $IM=M$ implies $im=m$.)
\end{lemma}
\begin{proof}
    Apply the previous corollary (\Cref{cor:cayley_hamilton}) to $\phi = \id_M$.
    Then, we have that
    $$(\id)^n + c_{n-1}(\id)^{n-1} + \dots + c_1(\id) + c_0 = 0,$$
    for some $c_0,\dots,c_{n-1}\in I$. This equation simplifies to
   $$(1+c_{n-1} +\dots+c_1+c_0)\id = 0.$$
   Since $c_{n-1},c_{n-2},\dots,c_0\in I$, we may take
   $i = -(c_0+c_1+\dots+c_{n-1})$.
\end{proof}
\begin{exercise}
Come up with a counterexample to show that 
the above lemma is not true if $M$ 
is not required to be finitely-generated.
\end{exercise}

The most prominent application of Nakayama's lemma
is in local ring.
Recall that a ring is local iff it has a unique maximal ideal.
If $R$ is local and $y\notin\mathfrak m$,
then $(y)=R$, so $y$ is a unit.
Thus, \textbf{if $R$ is a local ring with maximal ideal $\mm$,
any element not in $\mm$ is a unit.}
\begin{corollary}[Nakayama on local ring]
    \label{cor:nakayama_local_ring}
    \index{Nakayama's lemma}
    If $M$ is finitely-generated $R$-module, where
    $R$ is a local ring.
    Suppose that $M = \mathfrak m M$, then $M=0$.
\end{corollary}
\begin{proof}
    Apply Nakayama again.
    There exists $a\in\mathfrak m$ such that $(1+a)M = 0$.
    However, $1+a$ is a unit! Hence, we are done.
\end{proof}
Let us rephrase the above corollary in to a more useful form.
\begin{corollary}[Nakayama on generators]
    \label{cor:nakayama_generation}
    \index{Nakayama's lemma}
    Let $R$ be a local ring, and $M$ be a finitely-generated
    $R$-module. 
    Suppose that $(m_1,\dots,m_n)\in M$
    such that the quotient projection 
    $(\overline m_1, \dots, \overline m_n)
    \in M/\mathfrak mM$ span $M/\mathfrak mM$.
    Then, $(m_1,\dots,m_n)$ generate $M$.
\end{corollary}
\begin{proof}
    Let $N$ be submodule generated by $m_1,\dots,m_n$.
    From the generation assumption,
    we know that $N+\mathfrak mM = M$,
    so by quotienting with $N$, we have
    $\mathfrak m(M/N) = M/N$, so $M/N$ is zero.
\end{proof}
This corollary is actually very powerful.
$M/\mathfrak mM$ is a $R/\mathfrak m$-module,
but $R/\mathfrak m$ is a field,
so $M/\mathfrak mM$ is actually a vector space,
and it's then easy to check whether the generator spans.
This reduces the question of generation
into the world of linear algebra.
%
%
%
%
\section{Exactness}
\subsection{Exact Sequences}
Exact sequences is an expressive language that allows
us to concisely and graphically express many statements about
maps between modules.
They are used more extensively in algebraic topology
to manipulate a sheer number of spaces and maps,
but for us, we will mostly use this as a language
to state theorems.
\begin{definition}[Exact sequences]
    The sequence of mappings
    $$L\stackrel f\longto M \stackrel g\longto N$$
    is \vocab{exact} iff $\Im f = \Ker g$.

    Similarly, the long sequence
    $$\dots\longto M_{i-1}\stackrel{f_{i-1}}\longto M_i
    \stackrel{f_i}\longto M_{i+1}\to\dots$$
    is exact at $M_i$ iff $\Im f_{i-1} = \Ker f_i$.
    The whole diagram is an \vocab{exact sequence}
    if it is exact at all $M_i$'s.
\end{definition}
Now, let's give some examples.
\begin{example}
We have some basic examples of exactness.
\begin{itemize}
    \item $0\longto L\stackrel\psi\longto M$ is exact iff $\psi$
	is injective.
    \item $M\stackrel\phi\longto N\longto 0$
	is exact iff $\phi$ is surjective.
    \item $0\longto L\stackrel\psi\longto M
	\stackrel\phi\longto N$ is exact iff $L\simeq \Ker\phi$.
    \item $L\stackrel\psi\longto M
	\stackrel\phi\longto N\longto 0$ is exact iff
	$\Coker\psi \simeq N$.
\end{itemize}
\end{example}
\begin{definition}[Short exact sequence]
A \vocab{short exact sequence} is an exact sequence of the form
$0\longto L\longto M\longto N\longto 0$,
\end{definition}
\begin{example}[Example of short exact sequences]
    \begin{itemize}
    \item If $R=\mathbb Z$, then we have
    $$0 \longto \ZZ \stackrel{\cdot 2}\longto\ZZ 
    \longto \ZZ/2\ZZ \longto 0.$$
    \item 
    More generally, if $N$ is a submodule of $M$,
    we have a short exact sequence
    $$0\longto N\longto M\longto M/N\longto 0.$$
    \item 
    Even more generally, if $\phi : M\to N$
    is a module morphism, then we have a short exact sequence
    $$0\longto \Ker\phi \longto M
    \stackrel\phi\longto \Im\phi\longto 0.$$
    (The map $\Ker\phi\longto M$ is inclusion map).
    \item For any modules $M,N$,
    we have the short exact sequence
    $$0\to M\to M\oplus N\to N\to 0,$$
    where the former is by $m\mapsto (m,0)$
    and the latter is by $(m,n)\mapsto n$.
    \end{itemize}
\end{example}
One source of exact sequence that we will use
very late in the text is \vocab{snake lemma}.
\begin{lemma}[Snake Lemma]
\label{lem:snake}
Suppose that there is a commutative diagram 
whose both rows are exact:
\begin{center}
\begin{tikzcd}
    0 \arrow[r]
    & L \arrow[d,"\alpha"] \arrow[r, "f"]
    & M \arrow[d, "\beta"] \arrow[r, "g"]
    & N \arrow[d, "\gamma"] \arrow[r]
    & 0 \\
    0 \arrow[r] 
    & L' \arrow[r, "f'", swap]
    & M' \arrow[r, "g'", swap]
    & N'  \arrow[r]
    & 0 
\end{tikzcd}
\end{center}
Then, there is a long exact sequence 
$$0 \longto \Ker\alpha \longto \Ker\beta 
\longto \Ker\gamma \stackrel\delta\longto \Coker\alpha 
\longto \Coker\beta \longto \Coker\gamma\longto 0.$$
\end{lemma}
\begin{proof}[Proof Sketch]
The proof of this is a long, arduous diagram chasing
with no insights required,
so you do not need to understand the proof
unless you really want to.
In fact, I will only tell the definition of all the maps
and leave the remaining checking to the (intersested) readers.

You should check yourself that all maps
$\Ker\alpha\to\Ker\beta$, $\Ker\beta\to\Ker\gamma$,
$\Coker\alpha\to\Coker\beta$, and $\Coker\beta\to\Coker\gamma$
are induced from $f$, $g$, $f'$, and $g'$, respectively.
Thus, the most non-obvious map to define is 
$\delta : \Ker\gamma \to \Coker\alpha$.

Let $n\in\Ker\gamma\subseteq N$.
Then, since $g$ is surjective, pick \textbf{any}
$m\in M$ such that $g(m)=n$.
Now, since diagram commutes, $g'(\beta(m)) = \gamma(g(m))
= 0$, so $\beta(m)\in \Ker g' = \Im f'$,
so there exists unique $\ell\in L'$ such that 
$f'(\ell) = \beta(m)$. We now declare 
$\delta(n) = [\ell]$, the class of $\ell$.

We have to show that this does not depend on the choice of
$m$. Suppose that I have another choice,
say $m^*$. Then, $g(m-m^*) = 0$, 
so $m-m^*\in\Ker g = \Im f$,
and so there exists $x\in L$ such that $f(x)=m-m^*$.
Then, if the values of $\ell$ obtained from $m$ and $m^*$ are
are $\ell$ and $\ell^*$, respecitvely, then
$$f'(\ell - \ell^*) = \beta(m-m^*)
= \beta(f(x)) = f'(\alpha(x))
\implies \ell-\ell^* = \alpha(x)\in \Im\alpha,$$
so $\ell$ and $\ell^*$ are mapped to the same class
in $\Coker\alpha = L'/\Im\alpha$,
implying that this map does not depend on the choice of $m$.

We omit all checks of exactness.
\end{proof}
%
%
%
%
%
\subsection{Exact Functor}
In commutative algebra, the most common type of results about exact sequences
is of the form whether a functor preserves exact sequence or not.
We begin with a (not very important) definition.
\begin{definition}[Additive Functor]
    Given rings $R,S$.
    We say that a covariant functor $F:\Mod R\to \Mod S$ is
    an \vocab{additive functor} if the map $\Hom_R(A,B)\to \Hom_S(F(A),F(B))$
    is an abelian group homomorphism.

    Similarly, a contravariant functor $F:\Mod R\to\Mod S$
    is an additive functor if the map $\Hom_R(A,B)\to 
    \Hom_S(F(B),F(A))$ is an abelian group homomorphism.
\end{definition}
\begin{example}
    \label{ex:functors}
    Here are some example of additive functors.
    \begin{itemize}
    \item \textbf{The $\Hom_R(P,\bullet)$ functor.}
    \notation{$\Hom_R(P,\bullet)$}{a covariant $\Hom$ functor}
    Fix an $R$-module $P$, we have the covariant functor
    \begin{align*}
    \Hom_R(P,\bullet) : \Mod R & \to\Mod R \\
    M &\mapsto \Hom_R(P,M).
    \end{align*}
    To define the induced map, for any $\phi : M\to N$,
    we need to define a map $\phi_*: \Hom(P,M)\to \Hom(P,N)$.
    This map is simply compose by $\phi$: mapping
    $\pi$ to $\phi\circ \pi$
    (so we have diagram $P\longto M \stackrel\phi\longto N$).
    One  can check that it is additive.
    
    \item \textbf{The $\Hom_R(\bullet, P)$ functor.}
    \notation{$\Hom_R(\bullet, P)$}{a contravariant $\Hom$ functor}
    Similarly, we have the contravariant functor
    \begin{align*}
    \Hom_R(\bullet,P) : \Mod R & \to\Mod R \\
    M &\mapsto \Hom_R(M,P).
    \end{align*}
    To define the induced map, for any $\phi : M\to N$,
    we need to define a map $\phi^* : \Hom(N,P)\to \Hom(M,P)$.
    This map is simply compose by $\phi$: mapping
    $\pi$ to $\pi\circ \phi$
    (so we have diagram $M\stackrel\phi\longto N \longto P$).
    \item \textbf{The ``Mod $I$'' functor}
    For any ideal $I\subset R$, we have the covariant functor
    \begin{align*}
        F : \Mod R &\to \Mod{R/I} \\
        M &\mapsto M/IM.
    \end{align*}
    For each map $\phi : M\to N$, one can define
    $\phi_* : M/IM \to I/IN$ so that
    it commutes with quotient projection.
    This functor is additive.
    \end{itemize}
\end{example}
\begin{definition}[Left and right exact functors]
    An additive covariant functor is 
    \begin{itemize}
    \item \vocab{left-exact} if for any left-exact sequence
    $0\to A\to B\to C$,
    the sequence 
    $$0\to F(A)\to F(B)\to F(C)$$ is exact.
    \item \vocab{right-exact} if for any right-exact sequence
    $A\to B\to C\to 0$,
    the sequence $$F(A)\to F(B)\to F(C)\to 0$$ 
    is exact.
    \item \vocab{exact} if it is both 
    left-exact and right-exact.
    (In particular, it perserves injectivity,
    surjectivity, kernels, and cokernels.)
    \end{itemize}
    Similarly, an additive contravariant functor is
    \begin{itemize}
    \item \vocab{left-exact} if for any exact sequence
    $A\to B\to C\to 0$,
    the sequence $0\to F(C)\to F(B)\to F(A)$ is exact.
    \item \vocab{right-exact} if for any exact sequence
    $0\to A\to B\to C$,
    the sequence $F(C)\to F(B)\to F(A)\to 0$ is exact.
    \item \vocab{exact} if it is both 
    left-exact and right-exact.
    \end{itemize}
\end{definition}
\begin{example}[Trivial Examples of Exact Functors]
    The identity functor is exact.
    The functor $M\mapsto M\oplus M$ is exact.
\end{example}
However, being exact is generally quite a strong condition.
Let us now investigate functors from \Cref{ex:functors}.
\begin{proposition}
    \label{prop:left_right_exact}
    \begin{alphenum}
    \item The functor $\Hom_R(P,\bullet)$ is left-exact
    (i.e., exact sequence $0\to A\to B\to C$
    gives exact sequence $0\to \Hom(P,A)\to\Hom(P,B)\to\Hom(P,C)$).
    \item The functor $\Hom_R(\bullet,P)$ is left-exact
    (i.e., exact sequence $A\to B\to C\to 0$
    gives exact sequence $0\to \Hom(C,P)\to\Hom(B,P)\to\Hom(A,P)$).
    \item The mod-$I$ functor is right-exact
    (i.e., exact sequence $A\to B\to C\to 0$
    gives exact sequence $A/IA\to B/IB\to C/IC\to 0$).
    \end{alphenum}
\end{proposition}
\begin{proof}
    \begin{alphenum}
    \item Given $0\to A\stackrel\phi\to B\stackrel\pi\to C$
    is exact. We need to show that
    $$0\longto F(A)\stackrel{\phi_*}\longto F(B)
    \stackrel{\pi_*}\longto F(C)$$ 
    is exact.

    First, we show exactness at $F(A)$,
    i.e., $\phi_*$ is injective.
    Suppose $\phi_*(g) =\phi\circ g = 0$ for some $g$,
    then for all $p\in P$, $\phi(g(p)) = 0$.
    However, $\phi$ is injective, so $g(p)=0$ for all $p\in P$,
    so $g=0$.

    Now, we show exactness at $F(B)$, i.e., $\Im\phi_* = \Ker\pi_*$.
    First, note that $\pi_*(\phi_*(f)) = f\circ\phi\circ\pi = 0$, so $\Im\phi_*\subseteq \Ker\pi_*$.
    In other direction, let $g\in\Hom_R(P,B)$
    such that $\pi_*(g)=\pi\circ g = 0$. We want to show that
    $g=\phi\circ f=\phi_*(f)$ for some $f$.
    We have that for all $p\in P$, $\pi(g(p))=0$,
    so $g(p)\in\Ker\pi = \Im\phi$.
    With that in mind, we can define $f(p)$ to be
    the (unique) preimage of $g(p)$ w.r.t. $\phi$,
    so $\phi(f(p)) = g(p)$. It is easy to see that
    $f$ is a module morphism, and then we are done.
    \item Left as an exercise.
    \item You can try to prove this now, but this is a special case of \Cref{thm:tensor_is_right_exact}
    that we will prove later. \qedhere
    \end{alphenum}
\end{proof}
\begin{exercise}
Come up with counterexamples to show that
the three functors in \Cref{prop:left_right_exact}
are not exact.
\end{exercise}
It's worth noting that these proofs are
particularly long and tedious, but it's pretty routine.
At each step, there is only one thing you could possibly do.


One might wonder if exact functor also holds for long exact
sequence. It turns out that this is the case.
\begin{lemma}
    Suppose $F$ is an exact functor, then
    for any exact chain $L\stackrel f\longto M\stackrel g\longto N$,
    the induced map
    $F(L)\stackrel {F(f)}\longto F(M)\stackrel {F(g)}\longto F(N)$
    is exact.
\end{lemma}
\begin{proof}
    We have three exact sequences, each implying
    a new one by applying $F$:
    \begin{align} 
    \label{eq:exact_1}
	0\to \Ker f\to L\to \Im f \to 0 &\ \implies\ 
	F(\Ker f)\to F(L) \to F(\Im f) \to 0 \\
    \label{eq:exact_2}
	0\to \Ker g \to M \to \Im g\to 0 
	 &\ \implies\ 0\to F(\Ker g)\to F(M)\to F(\Im g)\to 0\\
        \label{eq:exact_3}
	0\to \Im g \to N \to \Coker g \to 0
	 &\ \implies\ 0\to F(\Im g)\to F(N).
    \end{align}
    The first exact sequence \eqref{eq:exact_1} implies that 
    $\Im F(f) = F(\Im f)$.
    On the other hand,
    \begin{align*}
    \Ker(F(g)) = \Ker(F(M)\to F(N))
    &\stackrel{\eqref{eq:exact_3}}= 
    \Ker(F(M)\to F(\Im g)) \\
    &\stackrel{\eqref{eq:exact_2}}= F(\Ker g) = F(\Im f).
    \end{align*}
    Combining these two information gives exactness at 
    $F(M)$.
\end{proof}
\section{Colimits}
\label{sec:colimit}
\emph{This section is only used in the proof of 
\Cref{chap:flatness}
(and only play supplementary role where it is used).
It may be skipped in the first reading.}

Several properties about module
(such as Cayley-Hamilton theorem and its consequences
in \Cref{sec:cayley}) requires hypothesis that
the module is finitely-generated.
Colimits will allow us to extend some of those facts
to infinitely-generated modules.

As an example, consider
$$\mathbb Z\left[\frac 12\right] = \left\{\frac{a}{2^n} : a,n\in\ZZ\right\}
\subset \mathbb Q.$$
This is not a finitely-generated $\ZZ$-module. 
However, one can write a chain of modules
$$\frac 12\mathbb Z \subset \frac 14\mathbb Z
\subset \frac 18\mathbb Z \subset\dots \subset
\mathbb Z\left[\frac 12\right],$$
and one has the inclusion map
$$\frac 12\mathbb Z \stackrel{\cdot 2}\longrightarrow
\frac 14\mathbb Z \stackrel{\cdot 2}\longrightarrow
\frac 18\mathbb Z \stackrel{\cdot 2}\longrightarrow \dots.$$
We will now generalize this a bit by introducing a
\vocab{directed system}.
Throughout this section, let $R$ be a ring.
\subsection{Directed System and Colimit}
\begin{definition}[Directed System]
Given a poset $I$, a \vocab{directed system} $M_I$
of $R$-modules for $I$ is given by
\begin{itemize}
    \item for each $i\in I$, we have an $R$-module $M_i$;
    \item if $i\leq j$, we have $f_{ij} : M_i\to M_j$
	such that if $i\leq j\leq k$, we have
	$f_{ik} = f_{jk} \circ f_{ij}$.
\end{itemize}
\end{definition}
\begin{example}
    If $I=\mathbb N$ ordered by the usual $\leq$,
    then a directed system on $I$ is a sequence of modules
    and morphism
    $$M_1 \stackrel{f_{12}}\longto M_2 \stackrel{f_{23}}\longto M_3 \longto \dots$$
\end{example}
Now, we want to take the limit of the directed system.
Given a directed system $M_I$, we want a universal $R$-module $M$
such that 
\begin{itemize}
    \item for each $i\in I$, we have $p_i : M_i\to M$.
    \item for each $i\leq j$, we have $p_j\circ f_{ij} = p_i$.
\end{itemize}
(i.e., we want all the arrows to be consistent).
\begin{center}
\begin{tikzcd}
    & & M \\
    M_{i} \arrow[rru, "p_i"] \arrow[r,"f_{ij}",swap]&
    M_{j} \arrow[ru, "p_j", swap] & & 
    \dots  \arrow[lu]
\end{tikzcd}
\end{center}
Therefore, we are going to state this
in terms of universal property.
\begin{definition}[Colimit]
    Given a directed system $M_I$ we say that an $R$-module
    $M$ is a \vocab{colimit} (or \vocab{direct limit})
    if 
    \begin{alphenum}
    \item there exists $p_i : M_i\to M$ such that
	whenever $i\leq j$, we have $p_j\circ f_{ij} = p_i$
	(i.e., everything commutes).
    \item for any $R$-module $N$ and mapping
	$q_i : M_i\to N$ satisfying the above condition,
	there exists unique $\phi : M\to N$
	such that $q_i = \phi \circ p_i$ for all $i\in I$.
    \begin{center}
    \begin{tikzcd}
        & & N \\
        & & M \arrow[u,"\exists!\phi",swap] \\
        M_{i} \arrow[rruu, "q_i", bend left] \arrow[rru, "p_i"] \arrow[r,"f_{ij}",swap]&
        M_{j} \arrow[ruu, "q_j"] 
        \arrow[ru, "p_j", swap] & & 
        \dots  \arrow[lu] \arrow[luu, bend right]
    \end{tikzcd}
    \end{center}
    \end{alphenum}
    and we denote it by $M := \varinjlim_{i\in I} M_i$.
    \notation{$\varinjlim_{i\in I} M_i$}{colimit}
\end{definition}
\begin{example}
    If the poset $I$ has no relation,
    then $\varinjlim M_i = \bigoplus_{i\in I} M_i$.
\end{example}
Let us prove that colimit exists and unique.
\begin{proposition}[Colimit exists and is unique]
    For each directed system $M_I$,  $\varinjlim M_i$ exists and is  unique
    (of course up to isomorphism).
\end{proposition}
\begin{proof}
    % \textbf{Uniqueness.}
    % This has the same pattern as all uniqueness proof
    % (e.g., in \Cref{subsec}
    % If $M$ and $M'$ works, then,
    % by applying UMP for $M$ with $N=M'$, we get
    % the mapping $\phi : M\to M'$.
    % Similarly, by applying UMP for $M'$ with $N=M$,
    % we get the mapping $\psi : M'\to M$.
    %
    % Now, $\phi\circ\psi = \id_{M'}$ by uniqueness
    % in the UMP for $M'$. Similarly, $\psi\circ\phi=\id_{M}$,
    % so $\phi$ and $\psi$ are inverses and act as
    % an isomorphism.
    The uniqueness part follows the same pattern
    as universal property arguments,
    e.g., in \Cref{sec:universal}.
    Thus, we focus on proving existence.
    We start with $\bigoplus_{i\in I} M_i$.
    Consider the submodule $K\subset\bigoplus_{i\in I}$
    generated by
    $$(0,\dots,0,m_i,\dots,-f_{ij}(m_i),0,\dots,0).$$
    We claim that $\left(\bigoplus_{i\in I}M_i\right) / K$ works
    as the colimit.
    The map $p_i$ is by the composition
    $M_i\to \bigoplus_{i\in I} M_i\to \left(\bigoplus_{i\in I} M_i\right)/K$.
    The commuting criterion is satisfied by
    the modding out by $K$.

    Now, given $N$, by universal properties for direct sum,
    we have the map $h : \bigoplus_{i\in I} M_i\to N$.
    Now, the condition that
    $q_i=q_j\circ f_{ij}$
    implies that $K\subset\Ker h$,
    so $h$ descends to $\left(\bigoplus_{i\in I} M_i\right)/K\to N$,
    and we are done.
\end{proof}
\begin{example}[A boring example]
    Consider the directed system with two elements,
    and consider the directed system $M_1\stackrel f\to M_2$.
    Then,
    $$\varinjlim M_I = \frac{M_1\oplus M_2}{(m_1, -f(m_1))}
     \simeq M_2$$
     More generally, if $I$ has a maximal element,
     then the colimit is going to be just
     the module at the maximum element.
\end{example}
\begin{exercise}
    \label{ex:rational}
    \begin{alphenum}
    \item Let $I=\mathbb N$ with relation being the usual
    $\leq$.
    Consider the diagram $M_i$ where $M_i=\ZZ$
    and the map $M_i\to M_{i+1}$ is multiplication by $2$.
    Prove that $\varinjlim M_i = \ZZ[1/2]$.
    \item Let $I=\mathbb N$ but with relation being
    divisibility (i.e., $i\prec j$ iff $i\mid j$).
    Let
    $$M_i = \frac 1i \mathbb Z\subset\mathbb Q.$$
    Prove that $\varinjlim M_i = \mathbb Q$.
    \end{alphenum}
\end{exercise}
At the beginning, we said that
the reason why we introduce colimits because
we want to study any modules from
finitely generated modules. Let's now make this more precise
by proving that colimits actually generate all modules.
\begin{proposition}[Colimit of finitely-generated
submodules]
    \label{prop:direct_limit_fingen}
    Given an $R$-module $M$, there exists a poset $I$
    and a directed system of finitely generated modules $M_I$
    such that $M = \varinjlim M_I$.
\end{proposition}
\begin{proof}
    This proof is mostly tautological.
    Let 
    $$I = \{\text{poset indexed by finitely generated submodules of }M\}.$$
    Element $i\in I$ correspond to module $M_i$.
    We let $i\leq j$ iff $M_i\subset M_j$,
    and the maps is by inclusion.
    Since $M_i\subset M$ and $\bigcup_{i\in I} M_i=M$, 
    we can easily verify from universal property $\varinjlim M_i = M$.
\end{proof}
\begin{remark}[Colimit indexed over a small category]
    One can think of $I$ as a category: the objects are
    elements of $I$, and there is a map $i\to j$
    iff $i\leq j$. In that case, a directed system $M_I$
    would be a functor from $I$ to $\Mod R$.
    
    Thus, one can generalize this by replacing a poset $I$
    with an arbitrary \vocab{small category} $C$
    (a small category is a category whose a 
    collection of objects and morphisms are both sets).
    A \vocab{directed system} indexed over $C$
    is a functor from $C$ to $\Mod R$,
    and the colimit can be defined very similarly.
    It turns out that colimit still exists and is unique.
\end{remark}
\subsection{Filtered Colimit Preserves Exactness}
It turns out that under a mild condition on $I$, colimit preserves exactness.
Let us make this more precise.
\begin{definition}
    A poset $I$ is \vocab{filtered} if and only if 
    any finite subset $S\subseteq I$ has an upper bound $u$
    such that $u\geq s$ for any $s\in S$.

    (Equivalently, $I\neq\emptyset$
    and any two elements in $I$ has an upper bound.)
\end{definition}
\begin{example}[Filtered posets]
    Most posets that we are considering the colimit
    is filtered.
    \begin{itemize}
    \item $\mathbb N$ with standard $\leq$ is filtered.
    \item $\mathbb N$ with divisibility relation
    (cf. \Cref{ex:rational})
    is filtered; given $m,n\in\NN$, $mn$
    is divisible by both $m$ and $n$, hence an upper bound.
    \item The poset corresponding to colimit in 
    \Cref{prop:direct_limit_fingen} is filtered.
    This is because for any finitely-generated submodules
    $N_1,\dots,N_k$, their sum $N_1+\dots+N_k$ is a
    finitely-generated submodule of $M$
    and is an upper bound of $\{N_1\dots,N_k\}$.
    \end{itemize}
    On the other hand, it is easy to come of with 
    poset that is not filtered:
    take $I=\{1,2,3\}$ with only relation $1\leq 2$ and $1\leq 3$.
\end{example}
\begin{definition}[Morphism of directed systems]
    Given directed system $M_I$ and $N_I$, a morphism from $M_I$
    to $N_I$ is given by $\phi_i : M_i\to N_i$ for each $i$
    such that for each $i\leq j$, the diagram
    \begin{center}
    \begin{tikzcd}
	M_i \arrow[r, "f_{ij}"]
	    \arrow[d, "\phi_i"]
	& M_j \arrow[d, "\phi_j"] \\
	N_i \arrow[r, "g_{ij}"]
	& N_j
    \end{tikzcd}
    \end{center}
    commutes. This promotes $\varinjlim$ to a functor.
\end{definition}
\begin{theorem}[Colimit preserves exactness]
    \label{thm:direct_lim_is_exact}
    Let $I$ be a \textbf{filtered} poset.
    Suppose that 
    $$L_I \stackrel{\{\psi_i\}}\longto
    M_I \stackrel{\{\phi_i\}}\longto N_I$$
    is a sequence of morphisms of directed systems
    such that
    $L_i \stackrel{\psi_i}\to M_i
    \stackrel{\phi_i}\to N_i$
    is exact for all $i\in I$.
    Then, we have the exact sequence
    $$\varinjlim L_I \longto \varinjlim M_I
    \longto \varinjlim N_I.$$
\end{theorem}
We start with the following lemma,
which is where the filtered condition comes into play.
\begin{lemma}
    \label{lem:direct_limit_exact}
    Let $I$ be a \textbf{filtered} poset,
    $M_I$ be a directed system,
    and $M = \varinjlim M_I$.
    Let $p_i : M_i\to M$ be the usual inclusion.
    Then,
    \begin{alphenum}
    \item $M = \bigcup_{i\in I} p_i(M_i)$
    \item Given $m_i$ such that $p_i(m_i) = 0$,
	there exists $j\geq i$ such that $p_{ij}(m_i)=0$.
    \end{alphenum}
\end{lemma}
\begin{proof}
    We do each part separately.
    \begin{alphenum}
    \item We use the construction of colimit that 
    $M$ is a quotient of $\bigoplus_{i\in I} M_i$.
    Given $m\in M$.
    Thus, 
    \begin{equation}
    \label{eq:colimit_image}
    m = p_{i_1}(m_{i_1}) + \dots + p_{i_n}(m_{i_n})
    \end{equation}
    for some $i_1,\dots,i_n\in I$ and $m_{i_1}\in M_{i_1}$,
    \dots, $m_{i_n}\in M_{i_n}$.
    
    Pick $j\geq i_1,i_2,\dots,i_n$.
    Thus, for any $\alpha\in\{1,\dots,n\}$,
    we have $p_{i_\alpha}(m_{i_\alpha})
    = p_j(f_{i_{\alpha}j}(m_{i_\alpha}))$
    by commuting condition,
    and so each of the summand in \eqref{eq:colimit_image} is in
    $p_j(M_j)$, implying the conclusion.
    \item Suppose that $p_i(m_i)=0$, then in $\bigoplus_{i\in I} M_i$, 
    we have that $m_i\in K$, where $K$ is the submodule
    in the explicit construction of colimits.
	In particular, there exists
	\begin{itemize}
	\item $i_1,i_2,\dots,i_n,j_1,j_2,\dots,j_n\in I$;
	\item $a_1,\dots,a_n\in R$; and
	\item $m_{i_1}\in M_{i_1}$, $m_{i_2}\in M_{i_2}$, 
	$\dots$, $m_{i_n}\in M_{i_n}$
	\end{itemize}
	such that
	\begin{equation}
    \label{eq:colimit}
    a_1(m_{i_1} - p_{i_1j_1}(m_{i_1}))
	+ \dots + a_n(m_{i_n} - p_{i_nj_n}(m_{i_n})) = m_i\quad \end{equation}
	Let $t$ be an upper bound of $i_1,\dots,i_n,j_1,\dots,j_n$.
	Thus, there is a map $M_{i_1},\dots,M_{i_n},M_{j_1},\dots,M_{j_n}\to
	M_t$. By universal property of direct sum,
	we get a map
	$$f : M_{i_1}\oplus \dots \oplus M_{i_n}\oplus M_{j_1}\oplus\dots\oplus
	M_{j_n}\to M_t.$$
	Taking $f$ on \eqref{eq:colimit}, we get that (now in $M_t$)
	$$a_1\big(p_{i_1t}(m_{i_1}) - p_{j_1t}(p_{i_1j_1}(m_{i_1}))\big)
	+ \dots
	+ a_n\big(p_{i_nt}(m_{i_n}) - p_{j_nt}(p_{i_1j_1}(m_{i_n}))\big)
	= p_{it}(m_i),$$
	and noticing that $p_{i_kt}(m_{i_k}) = p_{j_kt}(p_{i_kj_k}(m_{i_k}))$,
	we see that the left hand side is zero, meaning that
	$p_{it}(m_i)=0$, as desired. \qedhere
    \end{alphenum}
\end{proof}
\begin{proof}[Proof of \Cref{thm:direct_lim_is_exact}]
    We use the following notations.
    \begin{itemize}
    \item Let $f_{ij} : L_i\to L_j$,
    $g_{ij} : M_i\to M_j$, and $h_{ij} : N_i\to N_j$
    be the transition functions in the directed systems
    $L_I$, $M_I$, and $N_I$.
    \item Let $L = \varinjlim L_I$, $M=\varinjlim M_I$,
    and $N=\varinjlim N_I$.
    \item Let the resulting maps
    be $\psi : L\to M$ and $\phi : M\to N$.
    \end{itemize}
    It suffices to show exactness at $M$.

    We show that $\Ker\psi \subseteq \Im\phi$
    Let $m\in M$ such that $\psi(m)=0$.
    By \Cref{lem:direct_limit_exact} (a), 
    $m$ is image of $m_i\in M_i$ for some $i$.
    Then, since $\psi(m)=0$, we get that $\psi_i(m_i)=0$.
    Thus, by \Cref{lem:direct_limit_exact} (b), there exists
    $j\geq i$ such that $h_{ij}(\psi_i(m_i))=0$.
    Thus, $\psi_j(g_{ij}(m_i))=0$.
    
    By exactness at $M_j$, there exists $\ell_j\in L_j$
    such that $\psi_j(\ell_j) = g_{ij}(m_i)$.
    Taking $p_j$ both sides gives 
    $p_j(\psi_j(\ell_j)) = p_j(g_{ij}(m_i))$.
    However, the left hand side is $\psi(p_j(\ell_j))$,
    while the right hand side is $p_i(m_i) = m$,
    so $m$ lies in the image.

    The proof that $\Im\phi\subseteq\Ker\psi$
    is essentially the reverse of this argument.
\end{proof}
\begin{exercise}
Complete the proof of \Cref{thm:direct_lim_is_exact}.
\end{exercise}
%
%
%
%
\section{Tensor Product}
\label{sec:tensor}
Tensor product is a way to take \emph{product} of modules.
For general modules, the result is pretty strange,
so let us briefly describe what happens in a simpler 
case of vector spaces.

Given two $k$-vector spaces $V$ and $W$,
one construct the tensor product,
which is a vector space $V\otimes_k W$ by 
considering a big vector space generated by symbols 
$v\otimes w$ for each $v\in V$ and $w\in W$,
modulo the following relations:
\begin{itemize}
    \item $av\otimes w = v\otimes aw = a(v\otimes w)$
    for every $a\in k$ and $v,w\in V$; and
    \item $(v+v')\otimes w = (v\otimes w) + (v'\otimes w)$
    for every $v,v'\in V$ and $w\in W$.
    \item $v\otimes (w+w') = (v\otimes w) + (v\otimes w')$
    for every $v\in V$ and $w,w'\in W$.
\end{itemize}

If $\{e_i\}_{i\in I}$ is a basis of $V$
and $\{f_j\}_{j\in J}$ is a basis of $W$,
then $V\otimes W$ has basis $\{e_i\otimes f_j : i\in I, 
j\in J\}$.
(Proving that this set spans $V\otimes W$ is easy,
but proving that this set is linearly independent requires work.)
In particular, $\dim(V\otimes_k W) = (\dim V)(\dim W)$.

One can define tensor product of modules similarly
to above (i.e., tensor product $M\otimes_R N$ of $R$-modules 
$M$ and $N$ is generated by $m\otimes n$,
modulo the relations similar to above).
However, this definition makes it inconvenient to deduce
properties. Thus, we change tack by considering defining
tensor product by a universal property.
Through this, we see that tensor product 
tied closely with bilinear map. In particular,
there is a bijective correspondence between 
$$\{\text{module morphisms } M\otimes_R N\to P\}
\ \longleftrightarrow\ \{\text{bilinear maps } M\otimes N\to P\}.$$

This construction takes time to get used to,
but we will see various applications throughout the entire notes.
The most prominent application being the ability to 
``change the coefficients'' from e.g.,
$\RR$-algebra to $\CC$-algebra.
\subsection{Definition}
\paragraph{Bilinear maps}
\begin{definition}[Bilinear maps]
    Given modules $M,N,P$, the map
    $$\phi : M\times N\to P$$
    is \vocab{bilinear} if and only if $\phi$ is $R$-linear
    in each entry.
\end{definition}
Each map
$$f_\phi : M\to \Hom_R(N,P)$$
induces a bilinear map
\begin{align*}
    \phi : M\times N&\to P \\
	(m,n) &\mapsto f_\phi(m)(n).
\end{align*}
This establishes a natural isomorphism between
\begin{equation}
\label{eq:bilinear_hom}
\{\text{bilinear maps }M\times N\to P\}
\ \longleftrightarrow\ 
\Hom_R(M,\Hom_R(N,P)).
\end{equation}
%
%
%
\paragraph{Universal property}
Tensor product can be thought of as
a ``universal module'' represent those bilinear maps.
\begin{definition}[Tensor products]
    \notation{$M\otimes_R N$}{tensor product}
    Given $R$-modules $M,N$, the \vocab{tensor product}
    $M\otimes_R N$ is an $R$-module $T$
    with a bilinear map $\phi : M\times N \to T$
    such that for any $R$-module $P$ and a bilinear map
    $\psi : M\times N\to P$,
    there exists unique module homomorphism
    $\pi : T\to P$ such that $\psi = \pi\circ\phi$.
    \begin{center}
    \begin{tikzcd}
        M\times N \arrow[rr, "\text{bilinear}"]
        \arrow[rd]
        & & P \\
        & T \arrow[ru]
    \end{tikzcd}
    \end{center}
\end{definition}
As usual with universal properties,
$M\otimes_R N$ is unique up to isomorphism.
Thus, we need to prove that $M\otimes_R N$ exists.
\begin{proof}[Proof that $M\otimes_R N$ exists]
    Take an enormous free module
    $$F = \bigoplus_{(m,n)\in M\times N} Re_{(m,n)}.$$
    We have the obvious map $M\times N\to F$.
    Now, we need to mod out by something to force maps to be bilinear.
    We let $S$ be the submodule generated by
    \begin{align*}
    re_{(m,n)} - e_{(rm,n)} \\
    re_{(m,n)} - e_{(m,rn)} \\
    e_{(m+m',n)} - e_{(m,n)} - e_{(m',n)} \\
    e_{(m,n+n')} - e_{(m,n)} - e_{(m,n')}.
    \end{align*}
    (i.e., the things we want it to be zero).
    Then, one can compose $M\times N\to F\to F/S$
    to get $\phi : M\times N\to F/S$, which is bilinear.

    Now, for any $P$ and $\psi:M\times N\to P$ bilinear,
    define $\widetilde\pi : F\to P$ by $e_{(m,n)} = \psi(m,n)$.
    Notice that since $\psi$ is bilinear,
    $\widetilde\pi$ descends to $\pi : F/S \to P$,
    one can take $M\otimes_R N = F/S$.
\end{proof}
We also commonly denote $m\otimes n$ to be
the image of $(m,n)$ under $M\times N\to M\otimes_R N$.
\notation{$m\otimes n$}{pure tensor}
From the construction above, we get that
\begin{center}
\fbox{$M\otimes_R N$ is generated by \vocab{pure tensors}
$\{m\otimes n : m\in M, n\in N\}$}. 
\end{center}
%
%
%
%
%
\subsection{Properties of Tensor Product}
As we can see in the construction,
the construction of tensor products 
is very annoying to work with.
Thus, we will prove properties of tensor product
by using its universal properties instead.
\begin{proposition}[Basic properties of tensor products]
    \label{prop:tensor}
    Let $R$ be a ring and $M,N,P$ be modules.
    \begin{alphenum}
    \item $M\otimes_R N \simeq N\otimes_R M$
    (and the map is $m\otimes n\mapsto n\otimes m$).
    \item $M\otimes_R R\simeq M$ (and the map is 
    $m\otimes r\mapsto rm$).
    \item $(M\oplus N)\otimes_R P
    \simeq (M\otimes_R P) \oplus (N\otimes_R P)$
    (and the map is $(m+n)\otimes p \mapsto (m\otimes p)
    + (n\otimes p)$)
    \item $(M\otimes_R N)\otimes_R P
    \simeq M\otimes_R (N\otimes_R P)$
    (and the map is $(m\otimes n)\otimes p
    \mapsto m\otimes (n\otimes p)$).
    \item $M\otimes R/I \simeq M/IM$
    for all ideal $I\subseteq R$.
    \end{alphenum}
\end{proposition}
\begin{proof}
    The proofs below all follows the same pattern:
    we construct maps in both directions using 
    universal property. This is best illustrated through 
    the first few proofs.
    \begin{alphenum}
    \item We have the bilinear map
    \begin{align*}
	M\times N &\to N\otimes_R M \\
	(m,n) &\mapsto n\otimes m,
    \end{align*}
    so by universal property, we get a map
    \begin{align*}
	f : M\otimes_R N &\to N\otimes_R M \\
	m\otimes n &\mapsto n\otimes m.
    \end{align*}
    Similarly, one can construct
    $g : N\otimes_R M \to M\otimes_R N$
    which takes $n\otimes m$ to $m\otimes n$.
    Note that $f(g(m\otimes n)) = m\otimes n$,
    so $f\circ g$ and $\id$ coincide at pure tensors.
    Since $M\otimes_R N$ is generated by pure tensors,
    it follows that $f\circ g = \id$. Similarly, $g\circ f = \id$, so we are done.
    \item 
    We have the bilinear map
    \begin{align*}
	M\times R &\to M \\
	(m,r) &\mapsto rm,
    \end{align*}
    so it projects down to the map $f : M\otimes_R R\to M$
    sending $m\otimes r$ to $rm$.
    In another direction, one has the map $g : M\to M\otimes_R R$ sending $m\mapsto m\otimes 1$.

    To check that they are inverses, note that 
    $f(g(m)) = f(m\otimes 1) = m$, so $f\circ g=\id$.
    Similarly, 
    $$g(f(m\otimes r)) = g(rm) = rm\otimes 1 
    = r(m\otimes 1) = m\otimes r,$$
    so $g\circ f$ and $\id$ coincide at pure tensors,
    implying that $g\circ f = \id$.
    \item 
    Same pattern as before.
    We have the bilinear map
    \begin{align*}
	(M\oplus N)\otimes_R P &\to (M\otimes_R P)\oplus
	(N\otimes_R P) \\
	(m,n),p &\mapsto m\otimes p, n\otimes p,
    \end{align*}
    inducing a map $f : (M\oplus N)\otimes_R P
    \to (M\otimes_R P)\oplus (N\otimes_R P)$.

    Next, we construct the map from $(M\otimes_R P)
    \oplus (N\otimes_R P)\to (M\oplus N)\otimes_R P.$
    We note that the map
    \begin{align*}
    M\times P &\to (M\oplus N)\otimes_R P \\
    m,p &\mapsto (m,0)\otimes p
    \end{align*}
    is bilinear, so it descend down to the map 
    $M\otimes_R P \to (M\oplus N)\otimes_R P$
    that takes $m\otimes p$ to $(m,0)\otimes p$.
    Similarly, there is a map 
    $N\otimes_R P \to (M\oplus N) \otimes_R P$
    that takes $n\otimes p$ to $(0,n)\otimes p$.
    Thus, we get a map 
    \begin{align*}
    g : (M\otimes_R P) \oplus (N\otimes_R P) 
    &\mapsto (M\oplus N)\otimes_R P \\
    m\otimes p, n\otimes p' &\mapsto (m,0)\otimes p 
    + (0,n)\otimes p'.
    \end{align*}
    We leave the reader to check that $f$ and $g$
    are inverses.
    \item 
    One can follow the same argument as above.
    The other way is define a symmetric triple tensor product
    $M\otimes_R N\otimes_R P$ by the universal property
    that any trilinear map $M\times N\times P\to Q$
    factors through $M\otimes_R N\otimes_R P$.
    Then, show that $(M\otimes_R N) \otimes_R P 
    \simeq M\otimes_R N\otimes_R P$.
    \item 
    Define the map $M\times R/I\to M/IM$ by
    $(m,[a]) \mapsto [am]$.
    This gives a map $f : M\otimes_R R/I\to M/IM$.
    In other direction, consider the map $M\to M\otimes_R R/I$
    by sending $m\in M$ to $m\otimes [1]$.
    Notice that $im$ (for $i\in I$ and $m\in M$) is mapped to $im\otimes [1] = i(m\otimes [1]) = (m\otimes [i]) = 0$,
    so the map sends $IM$ to $0$, 
    hence descends down to $g : M/IM\to M\otimes R/I$.
    It's easy to see that $f$ and $g$ are inverses.
    \qedhere
    \end{alphenum}
\end{proof}
\begin{remark}
    Part (c) works with infinite direct sum:
    $M \otimes_R \bigoplus_{i\in I} N_i 
    = \bigoplus_{i\in I}(M\otimes_R N_i)$.
    However, it does not work for infinite products.
\end{remark}
Let's do some examples of tensor product.
\begin{example}[Computing tensor products]
    \begin{itemize}
    \item $R^{\oplus n}\otimes_R M = M^{\oplus n}$ (by \Cref{prop:tensor} (b) and (c)).
    \item In particular, $R^{\oplus m} \otimes_R R^{\oplus n}
    = R^{\oplus mn}$. 
    If $\{e_1,\dots,e_m\}$ and $\{f_1,\dots,f_n\}$
    are bases of $R^{\oplus m}$ and $R^{\oplus n}$, respectively,
    then $R^{\oplus m}\otimes_R R^{\oplus n}$
    has basis $\{e_i\otimes f_j : 1\leq i\leq m, 
    1\leq j\leq n\}$.
    A special case of this is when $R=k$,
    in which we recover tensor product of vector spaces.
    \item We have 
	$R/I\otimes R/J = R/(I+J)$ because by 
        \Cref{prop:tensor} (e),
	\begin{align*}
        R/I\otimes_R R/J &= (R/I)\ /\ (J(R/I)) \\
	&= (R/I)\ /\ ((I+J)/I) \\
        &= R/(I+J).
        \end{align*}
    \item In particular, 
        $\ZZ/m\ZZ \otimes_\ZZ \ZZ/n\ZZ \simeq 
        \ZZ/\gcd(m,n)\ZZ$.
        
    \item By PID structure theorem (\Cref{thm:pid_structure}), every finitely-generated $\ZZ$-module is a direct sum of 
        $\ZZ$ or $\ZZ/n\ZZ$. 
        Thus, we know how to compute tensor product of any two finitely-generated $\ZZ$-modules.
        For example,
        we can compute
    $$(\ZZ/6\ZZ \oplus \ZZ) \otimes_\ZZ
    (\ZZ/3\ZZ \oplus \ZZ/7\ZZ)$$
    by expanding four terms and computing the direct
    sum of these. We get that the tensor product is
    $$\ZZ/3\ZZ \oplus \ZZ/3\ZZ \oplus 0 \oplus \ZZ/7\ZZ.$$
    \end{itemize}
\end{example}
%
%
%
%
%
\subsection{Right Exactness of Tensor Product}
Let $R$ be a ring, and let $M$ be an $R$-module.
We have the functor
\begin{align*}
    \bullet \otimes_R M : \Mod R &\to \Mod R \\
    N &\mapsto N\otimes_R M.
\end{align*}
\notation{$\bullet\otimes_R M$}{the tensor functor}%
To define this functor, we need to, given 
morphism $\phi : N\to P$, define a map
$N\otimes_R M\to P\otimes_R M$. 
To do this, we note that the map 
\begin{align*}
    N\times M &\to P\otimes_R M \\
    (n,m) &\mapsto \phi(n)\otimes m
\end{align*}
is bilinear, and so by universal property of tensor,
we get the map $\phi_* : N\otimes_R M\to P\otimes_R M$
(the map is also often denoted $\phi\otimes\id$)
that sends $n\otimes m$ to $\phi(n)\otimes m$.

It turns out that this functor is right-exact.
\begin{theorem}[Tensor is Right Exact]
    \label{thm:tensor_is_right_exact}
    Fix an $R$-module $M$, the functor $M\otimes_R \bullet$
    is right-exact.
\end{theorem}
This theorem is useful in computing tensor product
of quotient because given the exact sequence
$$N \stackrel f\longto M\longto M/N \longto 0,$$
then the theorem gives the exact sequence
$$N\otimes_R P \stackrel{f_*}\longto M\otimes_R P\longto M/N
\otimes_R P\longto 0,$$
and so 
$$M/N\otimes_R P = \Coker f_*.$$
Therefore, given $M\otimes_R P$ and $N\otimes_R P$,
one can compute $M/N\otimes_R P$.

The idea of this proof is the following proposition.
\begin{proposition}[Tensor and Hom are adjoint]
    \label{prop:tensor_hom_adjunction}
    For any $R$-modules $M,N,P$, we have
    $$\Hom_R(M\otimes_R N, P) = \Hom_R(M,\Hom_R(N, P))$$
\end{proposition}
\begin{proof}
    By universal properties, there is an isomorphism 
    between 
    $$\Hom_R(M\otimes_R N, P)\ \longleftrightarrow\ 
    \{\text{bilinear maps } M\times N\to P\}.$$
    However, from \eqref{eq:bilinear_hom},
    this has a correspondence to 
    $\Hom_R(M, \Hom_R(N,P))$.
\end{proof}
\begin{lemma}[Converse of exactness of $\Hom$]
    \label{lem:hom_exactness}
    Let $R$ be a ring, and let $A$, $B$, $C$ be $R$-modules.
    Suppose that we have a sequence of maps
    $$A\stackrel f\longto B\stackrel g \longto C\longto 0.$$
    If the sequence after applying $\Hom_R(\bullet,T)$
    functor
    $$0\longto \Hom_R(C,T) \longto \Hom_R(B,T)\longto \Hom_R(A,T)$$
    is exact for all $R$-module $T$, 
    then the $A\to B\to C\to 0$ is exact.
\end{lemma}
\begin{proof}
    The proof is in three steps.
    \begin{enumerate}
    \item \textbf{$\Im f\subseteq \Ker g$.}
    Take $T=C$. Then, the identity map $\id\in\Hom_R(C,T)$
    is sent to $\id\circ f\in \Hom_R(B,T)$ 
    and then $\id\circ f\circ g \in\Hom_R(A,T)$,
    so by exactness at $\Hom_R(B,T)$, $\id\circ f\circ g = 0$,
    so $f\circ g = 0$, which implies 
    $\Im f\subseteq \Ker g$.
    \item \textbf{$g$ is surjective.}
    We pick $T=C/\Im g$ and let 
    $\pi : C\to C/T$ be the projection map.
    Then, note that under the map $\Hom_R(C,T) \to \Hom_R(B,T)$, element $\pi\in\Hom_R(C,T)$
    is sent to $\pi\circ g=0$.
    However, this map is injective by exactness at $\Hom_R(C,T)$, 
    so $\pi=0$, and so $g$ is surjective.
    \item \textbf{$\Ker g\subseteq\Im f$.}
    We set $T=B/\Im f$ and let $\phi : B\to T$ be 
    the projection map.
    Under the map $\Hom_R(B,T)\to \Hom_R(A,T)$,
    element $\phi\in\Hom_R(B,T)$ is sent to
    $\phi\circ f = 0$.
    Thus, $\phi\in \Ker(\Hom_R(B,T)\to\Hom_R(A,T))$.
    
    By exactness at $\Hom_R(B,T)$, $\phi$ is in image from $\Hom_R(C,T)$.
    In other words, there exists $\psi : C\to B/\Im(f)$
    such that $\psi\circ g = \phi$.
    Therefore, if $b\in \Ker g$, then 
    $g(b)=0$, so $\phi(b)=0$, implying that $b\in\Im f$.
    \qedhere
    \end{enumerate}
\end{proof}
Now, we are ready to prove \Cref{thm:tensor_is_right_exact}
\begin{proof}[Proof of \Cref{thm:tensor_is_right_exact}]
    Given the exact
    $$L\to N\to P \to 0,$$
    we want to show that
    $$M\otimes_R L\to M\otimes_R N\to M\otimes_R P\to 0$$
    is exact. By \Cref{lem:hom_exactness}, it suffices to show that
    $$0\to \Hom_R(M\otimes_R P, T)
    \to \Hom_R(M\otimes_R N, T)
    \to \Hom_R(M\otimes_R L, T)$$
    is exact. However, by Tensor-Hom adjunction
    (\Cref{prop:tensor_hom_adjunction}),
    we have $\Hom_R(M\otimes_R P, T)\simeq \Hom_R(P,\Hom_R(M,T))$.
    Thus, it suffices to show that
    $$0\to \Hom_R(P,\Hom_R(M,T))
    \to \Hom_R(N,\Hom_R(M,T))
    \to \Hom_R(L, \Hom_R(M,T))$$
    is exact, but this follows from $\Hom_R(\bullet,
    \Hom_R(M,T))$ is left-exact by \Cref{prop:left_right_exact} (a).
\end{proof}
\begin{remark}
    A special case of \Cref{thm:tensor_is_right_exact} 
    is when $M=R/I$, in which case
    $N\mapsto N\otimes_R R/I = N/IN$,
    giving the mod-$I$ functor.
    This proves \Cref{prop:left_right_exact} (c).
\end{remark}
\paragraph{Flat Modules}
\begin{definition}[Flat modules] 
\label{def: flat}
    A module $M$ is \vocab[flat module]{flat} if $M\otimes_R \bullet$ is exact
    (equivalent to preserve injection).
\end{definition}
\begin{example}[Examples of flat modules]
    \label{ex:flat}
    \begin{itemize}
    \item Free modules $R^{\oplus n}$ are flat.
    \item  More generally, $(M_i)_{i\in I}$ is flat
    for all $i$ if and only if  $\bigoplus_{i\in I} M_i$ is flat.
    \item $\ZZ/2\ZZ$-module is 
    not a flat $\ZZ$-module because 
    tensoring $\mathbb Z \stackrel{\cdot 2}{\longrightarrow} \mathbb Z$ with $\mathbb Z/2\mathbb Z$
    gives a zero map.
    \item More generally, any flat module is torsion-free:
    if $M$ is a flat $R$-module,
    then there is no $a\in R$ that is not a zero divisor
    such that $am=0$ for all $m\in M$.
    \item In particular, using the above bulllet point when $R=\ZZ$ and applying PID Structure theorem (\Cref{thm:pid_structure}), we deduce that every \textbf{finitely-generated} flat $\ZZ$-module is free (i.e. of the form $\ZZ^{\oplus n}$).
    \item However, we will prove later in \Cref{cor:localization_flat}
    that $\QQ$ is a flat $\ZZ$-module.
    \end{itemize}
\end{example}

We will study flat modules in detail in 
\Cref{chap:flatness}.

\subsection{Base Change}
\begin{definition}[$R$-algebra]
    Given a ring $R$, an $R$-\vocab{algebra} is a ring $S$
    with homomorphism $f : R\to S$.
    It has the obvious $R$-module structure.
\end{definition}
Let $S$ be an $R$-algebra.
Recall the ``restriction of scalar'' functor that given
an $S$-module $N$, one can construct the $R$-module
$_R N$ by forgetting multiplication in $S$.
What about other direction?

In another direction, we have the following construction.
\begin{definition}[Base Change]
Given an $R$-module $M$,
the \vocab{base change} of $M$ to an $S$-module 
is the module $M\otimes_R S$.
It has the structure of $S$-module,
and the scalar multiplication is defined by
$$x(m\otimes y) = m\otimes xy.$$
\end{definition}
\begin{example}
If $M=\mathbb R^3$, $R=\mathbb R$, and $S=\mathbb C$,
then $M\otimes_R S = \mathbb C^3$.
This operation is not inverse with restriction of scalars,
because if we restrict $\mathbb C^3$ back to $\RR$-module, we will get $\mathbb R^6$ instead.
\end{example}

If $M$ is an $R$-algebra, it turns out that you
get the algebra structure as well.
\begin{proposition}[Tensor product of algebras]
    Let $R$ be a ring and $S,T$ be $R$-algebras,
    then $S\otimes_R T$ is an $R$-algebra,
    and the multiplication structure is given by 
    $$(a\otimes b)(a'\otimes b') = aa'\otimes bb'$$
\end{proposition}
\begin{proof}
    To promote an $R$-module to $R$-algebra,
    we need to define the multiplication map 
    between two elements in $S\otimes_R T$.
    We have the map
    \begin{align*}
	(S\times T) \times (S\times T) &\to S\otimes_R T \\
	a, b, a', b' &\mapsto aa'\otimes bb',
    \end{align*}
    which is bilinear in $a$ and $b$ (when $a'$ and $b'$ are fixed) 
    and bilinear in $a'$ and $b'$ (when $a$ and $b$ are fixed).
    Thus, by universal property of tensor product,
    it descends down to the desired multiplication map.
    One can check that this satisfies all ring axioms
    (the identity is $1\otimes 1$, and associativity,
    commutativity, and distributivity holds).
\end{proof}

To compute base change, we have the following useful lemmas.
\begin{lemma}[Base change polynomial rings]
\label{lem:base_change}
Let $S$ be an $R$-algebra. Then,
\begin{alphenum}
    \item We have $R[x_1,\dots,x_n]\otimes_R S
    \simeq S[x_1,\dots,x_n]$ as $S$-algebras.
    \item Suppose that $f_1,\dots,f_m\in R[x_1,\dots,x_n]$
    and $g_i$ is the image of $f_i$ in $S[x_1,\dots,x_n]$, then
    $$\frac{R[x_1,\dots,x_n]}{(f_1,\dots,f_m)}
    \otimes_R S = \frac{S[x_1,\dots,x_n]}{(g_1,\dots,g_m)}$$
    as $S$-algebras.
\end{alphenum}
\end{lemma}
\begin{proof}
\begin{alphenum}
    \item Follows the same pattern as all tensor product proofs in \Cref{prop:tensor},
    so we will not go in full detail.
    In one direction, we can send
    $sx_1^{e_1}\cdots x_n^{e_n}$ to
    $s\otimes x_1^{e_1}\cdots x_n^{e_n}$.
    To find the inverse map, note that
    \begin{align*}
        R[x_1,\dots,x_n]\otimes_R S &\to S[x_1,\dots,x_n] \\
        rx_1^{e_1}\cdots x_n^{e_n}\otimes s &\mapsto
        rsx_1^{e_1}\cdots x_n^{e_n}
    \end{align*}
    is bilinear, giving the desired inverse map
    by universal property of tensor.
    \item We use right-exactness of tensor product
    \Cref{thm:tensor_is_right_exact}. Tensor
    $$(f_1,\dots,f_m) \longto R[x_1,\dots,x_n]
    \longto \frac{R[x_1,\dots,x_n]}{(f_1,\dots,f_m)}
    \longto 0$$
    with $S$ to get a right-exact sequence
    $$(g_1,\dots,g_m) \longto S[x_1,\dots,x_n]
    \longto \frac{R[x_1,\dots,x_n]}{(f_1,\dots,f_m)}
    \otimes_R S
    \longto 0,$$
    which implies the result. \qedhere
\end{alphenum}
\end{proof}
Thus, for example, 
$$\frac{\RR[x,y]}{(y^3-x^2-2025x)}
\otimes_\RR \CC 
= \frac{\CC[x,y]}{(y^3-x^2-2025x)}.$$
Therefore, base change is essentially a process
to change the coefficient rings to $\CC$.
\begin{example}
We have
\begin{align*}
    \CC \otimes_\RR \CC 
    &= \frac{\RR[x]}{(x^2+1)}
    \otimes_\RR \CC \\
    &= \frac{\CC[x]}{(x^2+1)} 
    \tag{\Cref{lem:base_change}}\\
    &= \frac{\CC[x]}{(x+i)} \times \frac{\CC[x]}{(x-i)} \\
    &= \CC \times \CC.
\end{align*}
Thus, $\CC\otimes_\RR \CC \simeq \CC\times \CC$ 
as $\RR$-algebra. Can you describe the bijection explicitly?
\end{example}

\section{Problems}
\begin{problem}
\label{prob:diff_xy}
Let $A=\mathbb C[x,y]/(xy)$ and $M$ be an $A$-module with generators $a$ and $b$ and single relation $xa+yb=0$. Prove that $M$ is not free.
\end{problem}
\begin{problem}
\label{prob:Zsqrt-5}
Let $A=\mathbb Z[\sqrt{-5}]$ and $I = (2, 1+\sqrt{-5})$ be an ideal. Prove that
\begin{alphenum}
\item $I\oplus I\simeq A\oplus A$ (as $A$-modules);
\item $I\otimes_A I\simeq A$.
\end{alphenum}
See \Cref{prob:dedekind_ideal} for a general result.

(Recommended approach for (b): right-exactness.)
\end{problem}
\begin{problem}
Let $A$ be a ring and $I$ be a finitely-generated ideal 
such that $I^2=I$. Prove that there exists $x\in I$
such that $I=(x)$ and $x^2=x$.
\end{problem}
\begin{problem}
\label{prob:projective_module_local_ring}
Let $A$ be a local ring, and let $P$ and $Q$ be $A$-modules such that $P\oplus Q \simeq A^{\oplus n}$ for some integer $n\geq 0$. Prove that $P$ and $Q$ are free.
\end{problem}
\begin{problem}
\label{prob:surjective_map}
Let $A$ be a ring and $M$ be a finitely-generated module. Prove that any surjective map $\phi : M\to M$ is an isomorphism.

(Hint: apply Nakayama's lemma on an $A[x]$-module.)
\end{problem}
\begin{problem}
\label{prob:projective}
Let $R$ be a ring, and let $P$ be an $R$-module.
\begin{alphenum}
\item Prove that the following are equivalent:
\begin{enumerate}[label=(\roman*)]
\item For every surjective map $\phi : P\to N$,
there exists a map $\psi : N\to P$ such that $\phi\circ \psi = \id$.
\item If $0\to M\to P\to N\to 0$ is an exact sequence,
then $P\simeq M\oplus N$.
\item There exists a module $Q$ such that $P\oplus Q$ is free.
\item For every surjective map of modules $\phi : M\to N$,
the induced map $\phi_* : \Hom_R(P,M)\to\Hom_R(P,N)$ 
is surjective (i.e., $\Hom_R(P,\bullet)$ is an exact functor).
\end{enumerate}
Module $P$ is said to be \vocab{projective} if it satisfies the above (equivalent) conditions.
Free modules are projective.
By \Cref{prob:projective_module_local_ring},
any finitely-generated projective modules over local ring is free.
It is true, but much more difficult to prove, that any projective module over local ring is free.
See \cite[Tag 0593]{stacks} for the proof.
\item Prove that projective modules are flat.%
\footnote{Conversely, you will see in \Cref{cor:flat_locally_free} that \textbf{finitely-generated} flat modules are projective.}
\item Prove that $\QQ$ is not a projective $\ZZ$-module.
\end{alphenum}
\end{problem}
\begin{problem}
\label{prob:colimit_tensor}
Prove that colimit commutes with tensor product.
In other words, given that $I$ is a filtered poset,
$(M_i)_{i\in I}$ is a directed system of $R$-modules,
and $N$ is an $R$-module, prove that
$$\varinjlim_{i\in I} (M_i\otimes_R N) 
\simeq \left(\varinjlim_{i\in I} M_i\right) \otimes_R N.$$

(Recommended approach: universal property.)
\end{problem}
\begin{problem}
Compute the following tensor products.
\begin{enumerate}[label=(\alph*)]
\item $k[x,y] \otimes_{k[x]} k[x,y]$, as a $k[x]$-algebra.
\item $\mathbb Q(\sqrt[3] 2) \otimes_{\mathbb Q} \mathbb Q(\sqrt[3] 2)$, as a $\QQ$-algebra 
(for this one, answer as a product of fields).
\item $\mathbb Q/\mathbb Z \otimes_{\mathbb Z} \mathbb Q/\mathbb Z$ as a $\ZZ$-module.
\end{enumerate}
\end{problem}
\begin{problem}
A rectangle is \emph{great} if and only if it has at least one rational side length. Suppose that a rectangle $R$ is tiled (with no overlap or empty space) by several axis-aligned great rectangles. Prove that $R$ is great.
\end{problem}